\SourceLeaf{21}
\par
It is easy to see that the entire surface of the spheroid is covered in this way if $p$ runs from $0$ to $180^0$, and $q$ from $0$ to $360^0$. We will further have
\[
\begin{aligned}
\lambda&=-A\sin p, &\quad\lambda'&=0\\
\mu&=B\cos p.\cos q, &\quad\mu'&=-B\sin p.\sin q\\
\nu&=C\cos p.\sin q, &\quad\nu'&=C\sin p.\cos q
\end{aligned}
\]
\[
\begin{aligned}
\mu\nu'-\nu\mu'&=BC\cos p.\sin p=ABC\sin p.\frac{x}{AA}\\
\nu\lambda'-\lambda\nu'&=AC\sin p^2.\cos q=ABC\sin p.\frac{y}{BB}\\
\lambda\mu'-\mu\lambda'&=AB\sin p^2.\sin q=ABC\sin p.\frac{z}{CC}
\end{aligned}
\]
\par
Hence, since $\sin p$ is positive everywhere within the limits considered here, we must set
\[
\mathrm{d}s=\mathrm{d}p.\mathrm{d}q.ABC.\psi.\sin p
\]
\par
Applying these formulas to Theorem Two, we obtain the volume of the body or (setting the density $=1$) its mass
\[
=\int\!\int\mathrm{d}p.\mathrm{d}q.ABC.\cos p^2.\sin p
\]
or, integrating first with respect to $q$,
\[
=2\pi\int\mathrm{d}p.ABC.\cos p^2.\sin p=\tfrac12\pi ABC\int\mathrm{d}p.(\sin p+\sin3p)
\]
This integral is to be taken from $p=0$ to $p=180^0$. It gives $\tfrac43\pi ABC$, as is known from other considerations.

\Article{13}
\par
To determine the attraction that the spheroid exerts on any point, assuming that the attraction of each element is inversely proportional to the square of its distance from the attracted point, we have $fr=\frac{1}{rr}$, $Fr=-\frac{1}{r}$, $\Phi r=r$. Let the attraction of the entire spheroid in the direction parallel and opposite to the $x$-coordinate axis be $=X$, and set $X=ABC\xi$. Then by Theorem Three we have
\[
X=\int\!\int\mathrm{d}p.\mathrm{d}q\frac{BCx\sin p}{rA}=\int\!\int\mathrm{d}p.\mathrm{d}q\frac{BC\cos p.\sin p}{r}
\]
and therefore
